\section{Results}
% One subsection per RQ. Keep the pattern consistent across RQs:
% motivation, method, results, summary box. The box at the end of each RQ
% gives the reader a one-line answer that can be scanned quickly.
 
\subsection{RQ1: [State the question]}
\motivation 
one paragraph on why this RQ matters. Skip this paragraph
if the motivation was already covered in Section III.

\approach 
the procedure used for this RQ only.

\results
prose description of what was observed, with supporting tables
and figures. Report effect sizes, not only p-values. Discuss the result
in the text, not only in the caption.
\summarybox{Answer to RQ1}{
XXX
}
 
\subsection{RQ2: [State the question]}
\motivation 
one paragraph on why this RQ matters. Skip this paragraph
if the motivation was already covered in Section III.

\approach 
the procedure used for this RQ only.

\results
prose description of what was observed, with supporting tables
and figures. Report effect sizes, not only p-values. Discuss the result
in the text, not only in the caption.

\summarybox{Answer to RQ2}{
XXX
}
 
\subsection{RQ3: [State the question]}
\motivation 
one paragraph on why this RQ matters. Skip this paragraph
if the motivation was already covered in Section III.

\approach 
the procedure used for this RQ only.

\results
prose description of what was observed, with supporting tables
and figures. Report effect sizes, not only p-values. Discuss the result
in the text, not only in the caption.
\summarybox{Answer to RQ3}{
XXX
}

 \begin{lstlisting}[caption=``Change Return Type'' refactoring causing a compiler error, label=sni_compile_error]
- public void removeSdchEncoding(
+ @@@Before@@
+ public void removeSdchEncoding(
 \end{lstlisting}



 \clearpage